\begin{thebibliography}{99}

\bibitem{Borel}
A. Borel,
\emph{Linear Algebraic Groups}, second edition,
Graduate Texts in Mathematics 126, Springer, 1991.
Corollary 1.4, Theorem 4.8, \S7, and Proposition 14.2.
\url{https://doi.org/10.1007/978-1-4612-0941-6}.

\bibitem{Deligne}
P. Deligne,
Hodge cycles on abelian varieties (notes by J.~S. Milne),
in \emph{Hodge Cycles, Motives, and Shimura Varieties},
Lecture Notes in Mathematics 900, Springer, 1982, pp.~9--100.
Proposition 3.6; revised electronic text (2003), p.~27:
\url{https://www.jmilne.org/math/Documents/Deligne82_2003.pdf}.

\bibitem{Albanese}
J.-P. Demailly,
\emph{Complex Analytic and Differential Geometry},
online book, version of 21 June 2012.
Chapter II, Theorems 5.3 and 8.8; Chapter VI, \S9.2.
\url{https://people.math.harvard.edu/~demarco/Math274/Demailly_ComplexAnalyticDiffGeom.pdf}.

\bibitem{Fulton}
W. Fulton,
\emph{Intersection Theory}, second edition,
Ergebnisse der Mathematik und ihrer Grenzgebiete (3) 2,
Springer, 1998.
Chapters 2 and 6; \S\S16.1, 18.3, 19.1--19.2.
\url{https://doi.org/10.1007/978-1-4612-1700-8}.

\bibitem{FultonHarris}
W. Fulton and J. Harris,
\emph{Representation Theory: A First Course},
Graduate Texts in Mathematics 129, Springer, 1991.
Lectures 14, 18--20; in particular \S19.4 and \S\S20.1--20.2.
\url{https://doi.org/10.1007/978-1-4612-0979-9}.

\bibitem{vanGeemen}
B. van Geemen,
Kuga--Satake varieties and the Hodge conjecture,
in \emph{The Arithmetic and Geometry of Algebraic Cycles}
(Banff, 1998), NATO Science Series C 548,
Kluwer, 2000, pp.~51--82.
Propositions 5.9 and 6.3; preprint version:
\url{https://arxiv.org/abs/math/9903146}.

\bibitem{Grauert}
H. Grauert,
Ein Theorem der analytischen Garbentheorie und die Modulr\"aume
komplexer Strukturen,
\emph{Publ. Math. Inst. Hautes \'Etudes Sci.} \textbf{5} (1960), 5--64;
correction, \emph{ibid.} \textbf{16} (1963), 35--36.
\url{https://www.numdam.org/item/PMIHES_1960__5__5_0/};
\url{https://www.numdam.org/item/PMIHES_1963__16__35_0/}.

\bibitem{GrothendieckHilbert}
A. Grothendieck,
Techniques de construction et th\'eor\`emes d'existence en g\'eom\'etrie
alg\'ebrique IV: les sch\'emas de Hilbert,
\emph{S\'eminaire Bourbaki}, exp.~221 (1960/61),
collected edition, vol.~6, pp.~249--276.
Theorem 3.2, p.~260; Proposition 3.8, p.~265.
\url{https://www.numdam.org/item/SB_1960-1961__6__249_0/}.

\bibitem{Hironaka}
H. Hironaka,
Resolution of singularities of an algebraic variety over a field of
characteristic zero, I, II,
\emph{Ann. of Math.} (2) \textbf{79} (1964), 109--203, 205--326.
Part I, Main Theorems I and I$^*$, pp.~132--133.
\url{https://doi.org/10.2307/1970486};
\url{https://doi.org/10.2307/1970547}.

\bibitem{Huybrechts}
D. Huybrechts,
\emph{Lectures on K3 Surfaces},
Cambridge Studies in Advanced Mathematics 158,
Cambridge University Press, 2016.
Chapter and result locators in this paper refer to the author's
prepublication draft:
\url{https://www.math.uni-bonn.de/people/huybrech/K3Global.pdf}.

\bibitem{Ehresmann}
I. Kol\'a\v r, J. Slov\'ak, and P.~W. Michor,
\emph{Natural Operations in Differential Geometry},
Springer, 1993, Lemma 9.2, p.~77.
Corrected electronic edition:
\url{https://www.mat.univie.ac.at/~michor/kmsbookh.pdf}.

\bibitem{KugaSatake}
M. Kuga and I. Satake,
Abelian varieties attached to polarized K3-surfaces,
\emph{Math. Ann.} \textbf{169} (1967), 239--242.
\url{https://doi.org/10.1007/BF01399540}.

\bibitem{Milne}
J.~S. Milne,
\emph{Introduction to Shimura Varieties},
revised electronic version, 16 September 2017.
Theorems 6.7--6.8 and Corollary 6.9, pp.~72--73.
\url{https://www.jmilne.org/math/xnotes/svi.pdf}.

\bibitem{SerreGAGA}
J.-P. Serre,
G\'eom\'etrie alg\'ebrique et g\'eom\'etrie analytique,
\emph{Ann. Inst. Fourier} \textbf{6} (1956), 1--42.
\S12, Theorems 1--3, pp.~19--20; \S19, Propositions 13--15,
pp.~29--30.
\url{https://doi.org/10.5802/aif.59}.

\bibitem{StacksCartier}
The Stacks Project Authors,
\emph{The Stacks Project},
Lemma 42.29.5, tag 02TB (Cartier intersection and fundamental cycles).
\url{https://stacks.math.columbia.edu/tag/02TB}.

\bibitem{StacksCorrespondences}
The Stacks Project Authors,
\emph{The Stacks Project},
\S45.3, tag 0FFZ (correspondences), especially
Lemma 45.3.1 and Lemma 45.3.6.
\url{https://stacks.math.columbia.edu/tag/0FFZ}.

\bibitem{StacksFamilies}
The Stacks Project Authors,
\emph{The Stacks Project},
tags 01V4 and 01V8 (smoothness and the fiber criterion),
01VU (high powers of ample sheaves),
and 0B5U (ampleness and cohomology).
\url{https://stacks.math.columbia.edu/tag/01V4};
\url{https://stacks.math.columbia.edu/tag/01V8};
\url{https://stacks.math.columbia.edu/tag/01VU};
\url{https://stacks.math.columbia.edu/tag/0B5U}.

\bibitem{Voisin}
C. Voisin,
\emph{Hodge Theory and Complex Algebraic Geometry I},
translated by L. Schneps,
Cambridge Studies in Advanced Mathematics 76,
Cambridge University Press, 2002.
\S\S6.1--6.3 and 7.1--7.3.
\url{https://doi.org/10.1017/CBO9780511615344}.

\end{thebibliography}
